% !TeX program = xelatex
%
% Лекция-сопровождение к лабораторной работе № 3
% «Машинный перевод и оценка качества».
%
% Сборка:  xelatex mlg_lab3_mt_quality.tex   (дважды)
% Оформление и макросы — в mlg_preamble.tex, общем для всех лекций курса.
% Все числа посчитаны на учебном корпусе (160 сегментов, столбцы ru_ref и
% ru_mt) и закреплены тестами репозитория gd_learn.
%
\documentclass[aspectratio=169,12pt]{beamer}

% Общая преамбула лекций курса. Подключается после \documentclass:
%
%     \documentclass[aspectratio=169,12pt]{beamer}
%     % Общая преамбула лекций курса. Подключается после \documentclass:
%
%     \documentclass[aspectratio=169,12pt]{beamer}
%     % Общая преамбула лекций курса. Подключается после \documentclass:
%
%     \documentclass[aspectratio=169,12pt]{beamer}
%     \input{mlg_preamble}
%
% Здесь только оформление и макросы — ни \title, ни текста слайдов, чтобы
% пять презентаций курса выглядели одинаково и правились в одном месте.
%
\usepackage{fontspec}
\usepackage{polyglossia}
\setmainlanguage{russian}
\setotherlanguage{english}
\setmainfont{Calibri}
\setsansfont{Calibri}
\setmonofont{Consolas}[Scale=0.92]

\usepackage{graphicx}
\usepackage{booktabs}
\usepackage{array}

% ---------------------------------------------------------------- оформление
\definecolor{ink}{HTML}{1B2433}
\definecolor{navy}{HTML}{2B4C6F}
\definecolor{brick}{HTML}{B03A2B}
\definecolor{moss}{HTML}{3E7D5A}
\definecolor{plum}{HTML}{6B5B95}
\definecolor{ash}{HTML}{8B9299}
\definecolor{paper}{HTML}{ECEFF3}

\usetheme{default}
\usecolortheme{default}
\setbeamercolor{structure}{fg=navy}
\setbeamercolor{frametitle}{fg=navy}
\setbeamercolor{normal text}{fg=ink}
\setbeamercolor{itemize item}{fg=navy}
\setbeamercolor{itemize subitem}{fg=ash}
\setbeamerfont{frametitle}{size=\Large,series=\bfseries}
\setbeamerfont{framesubtitle}{size=\small,series=\mdseries}
\setbeamercolor{framesubtitle}{fg=ash}
\setbeamertemplate{navigation symbols}{}
\setbeamertemplate{itemize item}{\raisebox{0.15ex}{\tiny$\bullet$}}
\setbeamertemplate{itemize subitem}{\raisebox{0.15ex}{\tiny--}}
\setbeamersize{text margin left=1.1cm, text margin right=1.1cm}

\setbeamertemplate{footline}{%
  \hbox{\begin{beamercolorbox}[wd=\paperwidth,ht=2.6ex,dp=1.4ex,leftskip=1.1cm,%
        rightskip=1.1cm]{}%
    {\color{ash}\scriptsize Основы машинного обучения \quad·\quad
     Цифровая лингвистика и локализация}\hfill
    {\color{ash}\scriptsize\insertframenumber}%
  \end{beamercolorbox}}}

% Врезка: связь слайда с лабораторной или с работой переводчика.
% Ширина считается точно: colorbox добавляет \fboxsep с каждой стороны,
% поэтому parbox должен быть на 2\fboxsep уже текстового поля — иначе
% на каждом слайде вылезает overfull hbox.
\newcommand{\врезка}[1]{%
  \par\vspace{0.5ex}\noindent
  \colorbox{paper}{%
    \parbox{\dimexpr\textwidth-2\fboxsep\relax}{%
      \vspace{0.7ex}\small #1\vspace{0.7ex}}}\par}

% Колонка таблицы заданной доли ширины, без растягивания пробелов.
% Объявлять её нужно здесь: внутри frame beamer пересканирует тело кадра,
% и #1 в определении вызывает «Illegal parameter number».
\newcolumntype{Л}[1]{>{\raggedright\arraybackslash}p{#1\textwidth}}

% Ссылка на интерактивную страницу к слайду. Playground'ы считают те же
% числа на том же корпусе, что и лабораторные, поэтому на них можно
% ссылаться прямо с лекции, не оговаривая расхождений.
\hypersetup{colorlinks=true,urlcolor=moss,linkcolor=moss}
\newcommand{\playground}[3]{%
  \par\vspace{0.4ex}%
  {\footnotesize\color{ash}$\rightarrow$~\href{#1}{\bfseries #2} — #3\par}}

\newcommand{\кейс}[1]{{\color{brick}\small\textbf{#1}}}
\newcommand{\числ}[1]{{\color{navy}\bfseries #1}}

%
% Здесь только оформление и макросы — ни \title, ни текста слайдов, чтобы
% пять презентаций курса выглядели одинаково и правились в одном месте.
%
\usepackage{fontspec}
\usepackage{polyglossia}
\setmainlanguage{russian}
\setotherlanguage{english}
\setmainfont{Calibri}
\setsansfont{Calibri}
\setmonofont{Consolas}[Scale=0.92]

\usepackage{graphicx}
\usepackage{booktabs}
\usepackage{array}

% ---------------------------------------------------------------- оформление
\definecolor{ink}{HTML}{1B2433}
\definecolor{navy}{HTML}{2B4C6F}
\definecolor{brick}{HTML}{B03A2B}
\definecolor{moss}{HTML}{3E7D5A}
\definecolor{plum}{HTML}{6B5B95}
\definecolor{ash}{HTML}{8B9299}
\definecolor{paper}{HTML}{ECEFF3}

\usetheme{default}
\usecolortheme{default}
\setbeamercolor{structure}{fg=navy}
\setbeamercolor{frametitle}{fg=navy}
\setbeamercolor{normal text}{fg=ink}
\setbeamercolor{itemize item}{fg=navy}
\setbeamercolor{itemize subitem}{fg=ash}
\setbeamerfont{frametitle}{size=\Large,series=\bfseries}
\setbeamerfont{framesubtitle}{size=\small,series=\mdseries}
\setbeamercolor{framesubtitle}{fg=ash}
\setbeamertemplate{navigation symbols}{}
\setbeamertemplate{itemize item}{\raisebox{0.15ex}{\tiny$\bullet$}}
\setbeamertemplate{itemize subitem}{\raisebox{0.15ex}{\tiny--}}
\setbeamersize{text margin left=1.1cm, text margin right=1.1cm}

\setbeamertemplate{footline}{%
  \hbox{\begin{beamercolorbox}[wd=\paperwidth,ht=2.6ex,dp=1.4ex,leftskip=1.1cm,%
        rightskip=1.1cm]{}%
    {\color{ash}\scriptsize Основы машинного обучения \quad·\quad
     Цифровая лингвистика и локализация}\hfill
    {\color{ash}\scriptsize\insertframenumber}%
  \end{beamercolorbox}}}

% Врезка: связь слайда с лабораторной или с работой переводчика.
% Ширина считается точно: colorbox добавляет \fboxsep с каждой стороны,
% поэтому parbox должен быть на 2\fboxsep уже текстового поля — иначе
% на каждом слайде вылезает overfull hbox.
\newcommand{\врезка}[1]{%
  \par\vspace{0.5ex}\noindent
  \colorbox{paper}{%
    \parbox{\dimexpr\textwidth-2\fboxsep\relax}{%
      \vspace{0.7ex}\small #1\vspace{0.7ex}}}\par}

% Колонка таблицы заданной доли ширины, без растягивания пробелов.
% Объявлять её нужно здесь: внутри frame beamer пересканирует тело кадра,
% и #1 в определении вызывает «Illegal parameter number».
\newcolumntype{Л}[1]{>{\raggedright\arraybackslash}p{#1\textwidth}}

% Ссылка на интерактивную страницу к слайду. Playground'ы считают те же
% числа на том же корпусе, что и лабораторные, поэтому на них можно
% ссылаться прямо с лекции, не оговаривая расхождений.
\hypersetup{colorlinks=true,urlcolor=moss,linkcolor=moss}
\newcommand{\playground}[3]{%
  \par\vspace{0.4ex}%
  {\footnotesize\color{ash}$\rightarrow$~\href{#1}{\bfseries #2} — #3\par}}

\newcommand{\кейс}[1]{{\color{brick}\small\textbf{#1}}}
\newcommand{\числ}[1]{{\color{navy}\bfseries #1}}

%
% Здесь только оформление и макросы — ни \title, ни текста слайдов, чтобы
% пять презентаций курса выглядели одинаково и правились в одном месте.
%
\usepackage{fontspec}
\usepackage{polyglossia}
\setmainlanguage{russian}
\setotherlanguage{english}
\setmainfont{Calibri}
\setsansfont{Calibri}
\setmonofont{Consolas}[Scale=0.92]

\usepackage{graphicx}
\usepackage{booktabs}
\usepackage{array}

% ---------------------------------------------------------------- оформление
\definecolor{ink}{HTML}{1B2433}
\definecolor{navy}{HTML}{2B4C6F}
\definecolor{brick}{HTML}{B03A2B}
\definecolor{moss}{HTML}{3E7D5A}
\definecolor{plum}{HTML}{6B5B95}
\definecolor{ash}{HTML}{8B9299}
\definecolor{paper}{HTML}{ECEFF3}

\usetheme{default}
\usecolortheme{default}
\setbeamercolor{structure}{fg=navy}
\setbeamercolor{frametitle}{fg=navy}
\setbeamercolor{normal text}{fg=ink}
\setbeamercolor{itemize item}{fg=navy}
\setbeamercolor{itemize subitem}{fg=ash}
\setbeamerfont{frametitle}{size=\Large,series=\bfseries}
\setbeamerfont{framesubtitle}{size=\small,series=\mdseries}
\setbeamercolor{framesubtitle}{fg=ash}
\setbeamertemplate{navigation symbols}{}
\setbeamertemplate{itemize item}{\raisebox{0.15ex}{\tiny$\bullet$}}
\setbeamertemplate{itemize subitem}{\raisebox{0.15ex}{\tiny--}}
\setbeamersize{text margin left=1.1cm, text margin right=1.1cm}

\setbeamertemplate{footline}{%
  \hbox{\begin{beamercolorbox}[wd=\paperwidth,ht=2.6ex,dp=1.4ex,leftskip=1.1cm,%
        rightskip=1.1cm]{}%
    {\color{ash}\scriptsize Основы машинного обучения \quad·\quad
     Цифровая лингвистика и локализация}\hfill
    {\color{ash}\scriptsize\insertframenumber}%
  \end{beamercolorbox}}}

% Врезка: связь слайда с лабораторной или с работой переводчика.
% Ширина считается точно: colorbox добавляет \fboxsep с каждой стороны,
% поэтому parbox должен быть на 2\fboxsep уже текстового поля — иначе
% на каждом слайде вылезает overfull hbox.
\newcommand{\врезка}[1]{%
  \par\vspace{0.5ex}\noindent
  \colorbox{paper}{%
    \parbox{\dimexpr\textwidth-2\fboxsep\relax}{%
      \vspace{0.7ex}\small #1\vspace{0.7ex}}}\par}

% Колонка таблицы заданной доли ширины, без растягивания пробелов.
% Объявлять её нужно здесь: внутри frame beamer пересканирует тело кадра,
% и #1 в определении вызывает «Illegal parameter number».
\newcolumntype{Л}[1]{>{\raggedright\arraybackslash}p{#1\textwidth}}

% Ссылка на интерактивную страницу к слайду. Playground'ы считают те же
% числа на том же корпусе, что и лабораторные, поэтому на них можно
% ссылаться прямо с лекции, не оговаривая расхождений.
\hypersetup{colorlinks=true,urlcolor=moss,linkcolor=moss}
\newcommand{\playground}[3]{%
  \par\vspace{0.4ex}%
  {\footnotesize\color{ash}$\rightarrow$~\href{#1}{\bfseries #2} — #3\par}}

\newcommand{\кейс}[1]{{\color{brick}\small\textbf{#1}}}
\newcommand{\числ}[1]{{\color{navy}\bfseries #1}}


\title{Лабораторная работа № 3}
\subtitle{Машинный перевод и оценка качества}
\date{}
\author{}

\begin{document}

% =====================================================================
\begin{frame}[plain]
  \vspace{1.2cm}
  {\color{brick}\small\bfseries ЛАБОРАТОРНАЯ РАБОТА № 3}

  \vspace{0.7cm}
  {\color{navy}\fontsize{30}{34}\selectfont\bfseries Машинный перевод\\
   и оценка качества\par}

  \vspace{0.5cm}
  {\color{ink}\large Что метрика измеряет на самом деле\par}

  \vspace{1.0cm}
  {\color{ash} 160 сегментов машинного перевода рядом с эталонным.
   Главный вопрос работы — не «какая метрика лучше», а «что именно
   каждая из них видит и чего не видит никогда».\par}
\end{frame}

% =====================================================================
\begin{frame}{Постановка}
  У каждого сегмента есть исходник, эталонный перевод и машинный перевод.
  Задача: решить, какие сегменты отдавать на редактуру — не читая все 160.

  \vspace{0.9em}
  \begin{columns}[T,onlytextwidth]
    \column{0.47\textwidth}
      {\color{navy}\bfseries Два подхода}\\[0.4em]
      {\small
      {\bfseries Метрики} сравнивают машинный перевод с эталоном и выдают
      число: BLEU, chrF, TER.\par
      \vspace{0.5em}
      {\bfseries Формальные проверки} смотрят на сам сегмент: целы ли
      плейсхолдеры, на месте ли отрицание, не потеряно ли предложение.}
    \column{0.47\textwidth}
      {\color{brick}\bfseries Почему это не одно и то же}\\[0.4em]
      {\small Метрика отвечает на вопрос «насколько похоже на эталон».\par
      \vspace{0.5em}
      Проверка отвечает на вопрос «сломается ли продукт».\par
      \vspace{0.5em}
      Работа показывает, что эти два вопроса расходятся сильнее, чем
      кажется.}
  \end{columns}
\end{frame}

% =====================================================================
\begin{frame}{Как считаются метрики}
  \begin{columns}[T,onlytextwidth]
    \column{0.47\textwidth}
      {\color{navy}\bfseries BLEU}\\[0.4em]
      {\small Доля совпавших словесных n-грамм (обычно до 4-грамм) плюс
      штраф за слишком короткий перевод.\par
      \vspace{0.6em}
      Работает со словами целиком: «сохранить» и «сохранена» для него
      разные слова без всякого родства.}
    \column{0.47\textwidth}
      {\color{navy}\bfseries chrF}\\[0.4em]
      {\small То же самое, но по символьным n-граммам.\par
      \vspace{0.6em}
      Поэтому на русском chrF обычно выше и спокойнее: словоформы одного
      слова имеют общие куски, и метрика их засчитывает.}
  \end{columns}

  \vspace{1.0em}
  \врезка{\textbf{Обе метрики измеряют совпадение формы.} Ни одна из них
  не знает, что такое смысл, — и это не упрощение для учебных целей, это
  их устройство.}
\end{frame}

% =====================================================================
\begin{frame}{Средние по корпусу}
  \begin{columns}[T,onlytextwidth]
    \column{0.44\textwidth}
      {\small
      На 160 сегментах:\par
      \vspace{0.7em}
      BLEU {\color{navy}\fontsize{24}{26}\selectfont\bfseries 0,46}\par
      \vspace{0.5em}
      chrF {\color{navy}\fontsize{24}{26}\selectfont\bfseries 0,64}\par
      \vspace{0.7em}
      Медиана BLEU — 0,37.}
    \column{0.50\textwidth}
      {\color{brick}\bfseries И что это значит?}\\[0.4em]
      {\small Само по себе — ничего.\par
      \vspace{0.6em}
      «BLEU 0,46» нельзя перевести ни в «качество 46~\%», ни в «нужно
      править 54~\% сегментов». Число сравнимо только с другим числом,
      посчитанным на том же корпусе тем же способом.\par
      \vspace{0.6em}
      Абсолютное значение BLEU не означает вообще ничего. Это первое, что
      нужно усвоить про метрики.}
  \end{columns}
\end{frame}

% =====================================================================
\begin{frame}{Диверсия: ломаем перевод нарочно}
  \framesubtitle{Сегмент s002, интерфейс: «Удалить файл «\{name\}»?»}

  \vspace{0.3em}
  {\small
  \begin{tabular}{@{}l@{\hspace{1.6em}}r@{\hspace{1.2em}}r@{\hspace{1.6em}}l@{}}
    \toprule
    что сломали & BLEU & chrF & формальная проверка \\
    \midrule
    ничего не трогали & 1,00 & 1,00 & --- \\
    {\color{brick}сломан плейсхолдер} & \числ{1{,}00} & \числ{1{,}00}
      & {\color{moss}\bfseries сработала} \\
    «удалить» $\rightarrow$ «стереть» & 0,63 & 0,65 & --- \\
    переставлены два слова & 0,55 & 0,68 & --- \\
    \bottomrule
  \end{tabular}}

  \vspace{0.9em}
  \врезка{\textbf{Сломанный плейсхолдер метрики не заметили вовсе.}
  \texttt{\{name\}} превратился в \texttt{(name\}}: пользователь увидит в
  диалоге «(name\}» вместо имени файла, а BLEU остался 1,00. Безобидная же
  замена слова на синоним роняет BLEU почти вдвое.}

  \playground{https://konkin-nikita.ru/gd_learn/mt/}{метрики машинного перевода}%
  {эти кнопки ломают перевод, числа пересчитываются сразу.}
\end{frame}

% =====================================================================
\begin{frame}{Почему BLEU ведёт себя именно так}
  \begin{columns}[T,onlytextwidth]
    \column{0.47\textwidth}
      {\color{navy}\bfseries Плейсхолдер}\\[0.4em]
      {\small Токенизатор отбрасывает пунктуацию, и \texttt{\{name\}} с
      \texttt{(name\}} превращаются в одинаковые токены. Для метрики
      изменения не произошло.}
    \column{0.47\textwidth}
      {\color{navy}\bfseries Синоним}\\[0.4em]
      {\small Одно заменённое слово рушит сразу все n-граммы, в которые оно
      входило: униграмму, биграмму, триграмму. На коротком сегменте это
      половина всех совпадений.}
  \end{columns}

  \vspace{1.0em}
  \врезка{\textbf{Вывод для практики.} Тяжесть ошибки для метрики зависит
  от того, сколько n-грамм она задела, а не от того, во что она обойдётся
  пользователю. Эти две величины не связаны.}
\end{frame}

% =====================================================================
\begin{frame}{Худший случай: смысл на противоположный}
  \framesubtitle{Сегмент s009, интерфейс}
  \begin{columns}[T,onlytextwidth]
    \column{0.52\textwidth}
      {\small
      Эталон:\par
      {\color{moss}\bfseries «Пароль должен содержать не менее 8 символов»}\par
      \vspace{0.6em}
      Потеряно одно слово:\par
      {\color{brick}\bfseries «Пароль должен содержать менее 8 символов»}\par
      \vspace{0.8em}
      BLEU \числ{0{,}50} \quad chrF \числ{0{,}90}\par
      Формальная проверка: {\color{brick}\bfseries молчит}}
    \column{0.44\textwidth}
      \врезка{Требование безопасности вывернуто наизнанку. chrF — 0,90:
      по форме сегменты почти совпадают.\par
      \vspace{0.5em}
      Проверка на отрицание ищет «не» перед глаголом и не знает про
      «не менее».}
  \end{columns}

  \vspace{0.5em}
  {\small Сравните с предыдущим слайдом: замена «удалить» на «стереть»
  ничего не ломает и роняет BLEU до 0,63. Инверсия требования к паролю —
  до 0,50. Разница между катастрофой и мелочью для метрики почти
  неразличима.}
\end{frame}

% =====================================================================
\begin{frame}{Формальные проверки: что они ловят}
  \begin{columns}[T,onlytextwidth]
    \column{0.47\textwidth}
      {\small
      На 160 сегментах проверки сработали на \числ{8}.\par
      \vspace{0.7em}
      Из них плейсхолдеры сломаны ровно в \числ{3} сегментах: s002, s012,
      s028.\par
      \vspace{0.7em}
      Остальные — потерянное отрицание и потерянное предложение.}
    \column{0.47\textwidth}
      {\color{moss}\bfseries Сильная сторона}\\[0.4em]
      {\small Проверка отвечает «да» или «нет» и указывает на конкретное
      место. Её результат можно вставить в договор: «сегменты со
      сломанными плейсхолдерами не принимаются».\par
      \vspace{0.6em}
      С BLEU так нельзя: порога, отделяющего годное от негодного, не
      существует.}
  \end{columns}

  \vspace{0.8em}
  \врезка{Восемь сегментов из 160 — это пять процентов корпуса. Столько
  работы нужно, чтобы поймать все поломки, которые ломают продукт.}
\end{frame}

% =====================================================================
\begin{frame}{Сколько стоит отбирать сегменты по метрике}
  \begin{columns}[T,onlytextwidth]
    \column{0.46\textwidth}
      {\small
      \begin{tabular}{@{}Л{0.44}r@{}}
        \toprule
        средство отбора & на проверку \\
        \midrule
        формальные проверки & \числ{8} \\
        BLEU ниже медианы & \числ{78} \\
        \bottomrule
      \end{tabular}\par
      \vspace{0.6em}
      Из 160 сегментов. Почти \textbf{десятикратная} разница в объёме ручной
      работы.}
    \column{0.48\textwidth}
      {\small
      Отбор «всё, что ниже медианы BLEU» отправляет на редактуру
      \числ{49~\%} корпуса — и всё равно пропускает поломки.\par
      \vspace{0.7em}
      {\color{brick}\bfseries Худшее из двух миров:} вдесятеро больше
      работы и при этом не полное покрытие.\par
      \vspace{0.7em}
      Метрика не предназначена для отбора отдельных сегментов. Она
      предназначена для сравнения систем на большом корпусе.}
  \end{columns}
\end{frame}

% =====================================================================
\begin{frame}{Слепое пятно, названное по имени}
  \begin{columns}[T,onlytextwidth]
    \column{0.52\textwidth}
      {\small
      Сегмент {\color{navy}\bfseries s002} в корпусе:\par
      \vspace{0.5em}
      эталон \quad \texttt{Удалить файл «\{name\}»?}\par
      машина \quad {\color{brick}\texttt{Удалить файл "(имя\}"?}}\par
      \vspace{0.8em}
      BLEU \числ{0{,}63} — {\bfseries выше} медианы 0,37.\par
      \vspace{0.5em}
      Отбор по BLEU этот сегмент пропустит. Формальная проверка его
      ловит.}
    \column{0.44\textwidth}
      \врезка{Плейсхолдер сломан и вдобавок переведён: пользователь
      увидит в диалоге буквально «(имя\}» вместо имени файла.\par
      \vspace{0.5em}
      Для BLEU это сегмент выше среднего качества.}
  \end{columns}

  \playground{https://konkin-nikita.ru/gd_learn/mt/}{метрики машинного перевода}%
  {все слепые пятна корпуса — отдельной вкладкой.}
\end{frame}

% =====================================================================
\begin{frame}{Ложные срабатывания и морфология}
  Правило «потеряно отрицание» ищет частицу «не». На русском этого мало:
  отрицание бывает встроено в слово.

  \vspace{0.6em}
  \begin{columns}[T,onlytextwidth]
    \column{0.47\textwidth}
      {\color{brick}\bfseries Ложные тревоги}\\[0.4em]
      {\small «Невозможно…», «отсутствуют…» — отрицание на месте, но
      выражено морфологически. Правило считает его потерянным.\par
      \vspace{0.5em}
      В корпусе таких сегментов два: s010 и s031.}
    \column{0.47\textwidth}
      {\color{moss}\bfseries Настоящая ошибка}\\[0.4em]
      {\small s083 — потеряно целое предложение. После уточнения правила
      сегмент обязан остаться в списке.\par
      \vspace{0.5em}
      Это и есть критерий правильной правки: ложные тревоги ушли,
      настоящая ошибка осталась.}
  \end{columns}

  \vspace{0.9em}
  \врезка{\textbf{Задание работы.} Уточнить правило так, чтобы оно
  учитывало морфологическое отрицание. Проверяется не формулировка
  регулярного выражения, а то, что s010 и s031 ушли, а s083 остался.}
\end{frame}

% =====================================================================
\begin{frame}{Что должно быть в отчёте}
  \begin{columns}[T,onlytextwidth]
    \column{0.47\textwidth}
      {\color{moss}\bfseries Обязательно}\\[0.5em]
      {\small
      1. Таблица диверсий с разбором, почему метрики отреагировали именно
      так.\par
      \vspace{0.4em}
      2. Сравнение стоимости отбора: 8 против 78.\par
      \vspace{0.4em}
      3. Уточнённое правило отрицания и доказательство, что s083 остался
      помеченным.}
    \column{0.47\textwidth}
      {\color{brick}\bfseries Не принимается}\\[0.5em]
      {\small «chrF лучше BLEU для русского языка».\par
      \vspace{0.6em}
      Это верное утверждение, приведённое без опоры. Нужно показать на
      конкретных сегментах, где именно chrF ведёт себя разумнее — и где
      он точно так же слеп, как BLEU.}
  \end{columns}
\end{frame}

% =====================================================================
\begin{frame}{Что унести из работы № 3}
  \begin{columns}[T,onlytextwidth]
    \column{0.47\textwidth}
      {\color{navy}\bfseries Про метрики}\\[0.5em]
      {\small Метрика измеряет совпадение формы с эталоном. Стоимость
      ошибки для пользователя она не измеряет и не может.\par
      \vspace{0.6em}
      Абсолютное значение BLEU не значит ничего. Сравнивать можно только
      два числа с одного корпуса.}
    \column{0.47\textwidth}
      {\color{brick}\bfseries Про профессию}\\[0.5em]
      {\small Когда заказчик говорит «поднимем BLEU до 0,6», ваша работа —
      спросить, какие ошибки при этом останутся.\par
      \vspace{0.6em}
      Формальные проверки дёшевы, объяснимы и попадают в договор. Метрики
      — нет.}
  \end{columns}

  \vspace{1.0em}
  \playground{https://konkin-nikita.ru/gd_learn/mt/}{метрики машинного перевода}%
  {сломать перевод и увидеть реакцию метрик — за один клик.}
  \playground{https://konkin-nikita.ru/gd_learn/labels/}{данные и разметка}%
  {порог BLEU против разметки MQM и каппа разметчиков (к заданию 8).}
\end{frame}

\end{document}
